% !TeX TXS-program:compile = txs:///pdflatex/[-shell-escape]
% !TeX TXS-program:pdflatex = pdflatex -synctex=1 -interaction=nonstopmode --shell-escape %.tex
\section{Endian Escapades}\label{sec:endian_escapades}

\subsection{Texto - Litte-Endian}\label{subsec:litte_endian}

Cualquier valor que ocupe más de un byte, se divide en bytes individuales
para poder almacenarse en memoria, ya que la memoria se direcciona byte a byte.

La mayoría de las arquitecturas almacenan el dígito \textit{menos significativo} a la izquierda
y el \textit{más significativo} a la derecha. Un ejemplo sería el \( 0x1234 \), se
almacenaría como \( 34 \) \( 12 \). A este tipo de arquitecturas se les denomina Litte-Endian.
Esto de aquí, se realiza por las siguiente manera:

\begin{enumerate}
	\item Las mayoría de las operaciones aritméticas, se realiza partiendo del \textit{litte\_endian} y esto que los componentes electrónicos sea más natural.
	\item La dirección de un valor es la dirección de su byte menor, por lo que leer 1 byte o 2 bytes,
	      equivale a la misma dirección, simplemente se lee menos bytes.
\end{enumerate}

\begin{table}[htbp]
	\centering
	\label{tab:example_little_and_big_endian}
	\resizebox{\textwidth}{!}{%
		\begin{tabular}{lllll}
			\toprule
			size of value & decimal value & hex value          & big endian bytes (NOT x86) & little endian bytes (x86) \\
			\midrule
			8 (1 byte)    & 65            & 0x41               & 41                         & 41                        \\
			\midrule
			16 (2 bytes)  & 4660          & 0x1234             & 12 34                      & 34 12                     \\
			\midrule
			32 (4 bytes)  & 1145258561    & 0x44434241         & 44 43 42 41                & 41 42 43 44               \\
			\midrule
			64 (8 bytes)  & 1145258561    & 0x0000000044434241 & 00 00 00 00 44 43 42 41    & 41 42 43 44 00 00 00 00   \\
			\bottomrule
		\end{tabular}%
	}
	\caption{Ejemplo de almacenamiento en \textbf{\textit{Little-Endian}} y \textbf{\textit{Big-Endian}}}
\end{table}

\subsection{Texto - Memory Order and Register Values}\label{subsec:memory_order_and_register_values}

El punto en el que es más facil confundirse es la forma entre el orden de la
memoria y el valor que se imprime desde un registro.

Esto no importa en sí, puesto que \textbf{\textit{x86}} lo organiza y
maneja por nosotros.

\begin{table}[htbp]
	\begin{center}
		\begin{tabular}{cc}
			\toprule
			\textbf{memory address order: } & 41 42 43 44 45 46 47 48 \\
			\midrule
			\textbf{register hex order: }   & 48 47 46 45 44 43 42 41 \\
			\midrule
			\textbf{rax value: }            & 0x4847464544434241      \\
			\bottomrule
		\end{tabular}
	\end{center}
	\caption{Ejemplo de como se almacena}\label{tab:example_store_number_little_endian}
\end{table}

\subsection{Texto - Bytes, Words, and Friends}\label{subsec:bytes_words_and_friends}

En x86 se va a estar trabajando siempre con diferentes tamaños de \textit{bits},
\( 2,4,6\ u \ 8 \), y es bueno concerlo porque siempre vamos a estar
trabajando con ellos:

\begin{table}[htbp]
	\begin{center}
		\begin{tabular}[c]{|l|l|l|l|l|}
			\hline
			Name & Bits & Bytes & Partial \mintinline{ASM}{rax} Access & Memory Access                                              \\
			\hline
			byte & 8    & 1     & \mintinline{ASM}{mov al, [rdi]}      & \mintinline{ASM}{mov BYTE PTR [rdi], 0x11 }                \\
			\hline
			word & 16   & 2     & \mintinline{ASM}{mov ax, [rdi]}      & \mintinline{ASM}{mov WORD PTR [rdi], 0x1122 }              \\
			\hline
			byte & 32   & 4     & \mintinline{ASM}{mov eax, [rdi]}     & \mintinline{ASM}{mov DWORD PTR [rdi], 0x11223344 }         \\
			\hline
			byte & 64   & 8     & \mintinline{ASM}{mov rax, [rdi]}     & \mintinline{ASM}{mov QWORD PTR [rdi], 0x1122334455667788 } \\
			\hline
		\end{tabular}
	\end{center}
	\caption{Acceso a \( 2,4,6,8 \) bits.}\label{tab:access_common_bits}
\end{table}

\subsection{Sign Extension}\label{subsec:sign_extension}
\subsection{Litte-Endian Bytes}\label{subsec:litte_endian_bytes}
\subsection{Qword by Qword}\label{subsec:qword_by_qword}
\subsection{Dword by Dword}\label{subsec:dword_by_dword}
\subsection{Word by Word}\label{subsec:word_by_word}
\subsection{Bytes by Bytes}\label{subsec:bytes_by_bytes}
\subsection{Cracking a Struct}\label{subsec:cracking_a_struct}
\subsection{Scrambled Struct}\label{subsec:scrambled_struct}
